\section{Sharp local metric reduction for every boundary size}\label{sec:rooted-truncation}
The five-point cone argument above proves the unrestricted classification.
We now give a separate metric theorem that applies to any number of designated
points and any radius. It replaces metric-cone enumeration for a \emph{single
rooted local problem}; compatibility between different roots remains a
separate requirement.

\begin{theorem}[Sharp rooted metric representative]\label{thm:rooted-truncation}
Let $M$ be an integer metric with distinguished point $o$, and let $R\ge0$ be
an integer. For each point $i$ put
\[
 D_i=M(o,i),\qquad H_{ij}=D_i+D_j-M(i,j),\qquad T=2R+1.
\]
Set $D'_i=D_i$ when $D_i\le R$ and $D'_i=T$ otherwise. For distinct points set
$H'_{ij}=\min(H_{ij},T)$ and
\[
 M'(i,j)=D'_i+D'_j-H'_{ij},\qquad M'(i,i)=0.
\]
Then $M'$ is an integer metric on the same points, with root distances $D'_i$.
It preserves the clipped rooted signature
\[
 \bigl(\min(D_i,R+1),\ \min(H_{ij},2R+1)\bigr)_{i,j},
\]
and satisfies
\[
 \max_i M'(o,i)\le2R+1,\qquad \diam(M')\le4R+2.
\]
The transformation is idempotent. Both constants are best possible:
four points force the root-distance bound and five points force the diameter
bound.
\end{theorem}
\begin{proof}
The metric inequalities are equivalent to the relevant relations
\[
 0\le H_{ij}\le2\min(D_i,D_j),\quad
 H_{ij}<D_i+D_j\ (i\ne j),\quad
 H_{ij}+H_{jk}-H_{ik}\le2D_j.
\]
Symmetry is immediate. If both $D_i,D_j$ are at most $R$, their mutual
distance is unchanged. If both exceed $R$, their new distance is at least
$T$. In the mixed case their new distance is at least $T-D_j\ge R+1$,
where $D_j$ is the smaller root distance. Thus distinct points remain distinct.

The triangle inequality through $j$ becomes
$H'_{ij}+H'_{jk}-H'_{ik}\le2D'_j$. If $D_j>R$, the left side is at most
$2T=2D'_j$. If $D_j\le R$, its two incident $H$ values are unchanged. When
$H_{ik}\le T$ the old triangle inequality applies. Otherwise
\[
 H'_{ij}+H'_{jk}-H'_{ik}\le4D_j-T\le2D_j.
\]
The root distances, signature identities, bounds, and idempotence now follow
from the definitions.

For root-bound sharpness take points $o,i,j,k$ with
\[
 M(o,i)=M(o,k)=M(i,j)=M(j,k)=R+1,\quad
 M(o,j)=T,\quad M(i,k)=2R+2.
\]
These distances form a metric. Its products satisfy
$H_{ij}=H_{jk}=T$ and $H_{ik}=0$. Any representative with the same clipped
signature therefore has $2D'_j\ge T+T-0$, so $D'_j\ge T$.

For diameter sharpness take four cyclically ordered points and a root.
Give each root distance and each cyclic-edge distance the value $T$, and
each opposite-pair distance $2T$. This is a metric. Products are $T$ on
cyclic edges and zero on opposite pairs. For each cyclic point the preceding
triangle inequality forces root distance at least $T$. Each opposite pair
still has product zero, so its distance is at least $2T=4R+2$.
\end{proof}

\begin{theorem}[Preservation of the complete local frame]\label{thm:rooted-frame}
At radius $R$, the preceding transformation preserves every positive profile,
every designated point in the ball, sup-norm adjacency, and all signed
neighbor-score coefficients toward designated points, in normalized
coordinates.
\end{theorem}
\begin{proof}
At root-distance $1\le r\le R$, write
$f_o=r$ and $f_i=D_i+\epsilon_i$ with $-r\le\epsilon_i\le r$.
The positive-frame conditions become
\[
 D_i+\epsilon_i>0,\qquad \epsilon_i+\epsilon_j\ge-H_{ij},\qquad
 \exists i:\epsilon_i=-r,
\]
and, for every $i$, either $\epsilon_i=-r$ or some $j\ne i$ satisfies
$\epsilon_i+\epsilon_j=-H_{ij}$. Positivity is automatic whenever $D_i>R$.
When $H_{ij}>2R$, its inequality is automatic and cannot be tight. All other
quantities in these tests are unchanged, giving the required bijection.

A designated point $c$ lies in the ball exactly when $D_c\le R$. Its
normalized coordinates are $r=D_c$, $\epsilon_c=-D_c$, and
$\epsilon_i=D_c-H_{ci}$ for $i\ne c$. Since $H_{ci}\le2D_c\le2R$, these
coordinates are unchanged. Finally, keeping every offset keeps every
coordinate difference. This preserves sup-norm distances, adjacency, and
all signed coefficients in the local harmonic inequalities.
\end{proof}

For $k$ designated points, a direct finite search now needs at most
\[
 (R+1)^{k-1}(2R+2)^{\binom{k-1}{2}}
\]
rooted signatures before the metric conditions and symmetries are imposed.
This is a uniform finite reduction for all $k,R$, with no classification of
metric cones as an input. For $k=5,R=2$, direct generation gives 13,022
canonical metric types. The stronger experimental local model leaves 1,226
feasible types. These latter computations are discovery checks, not a
replacement for the independently certified all-root proof in the main
classification.

\paragraph{A necessary limitation.}
The theorem preserves the data at the chosen root only. At $R=2$, a
three-point metric with $D_a=3$, $D_b=4$, and $M(a,b)=1$ becomes
$D'_a=D'_b=5$, $M'(a,b)=5$. The root-$o$ signature is preserved while the
root-$a$ signature changes. Thus one cannot obtain a simultaneous global
cap merely by applying the transformation at different roots.

\paragraph{Formal and independent verification.}
\texttt{RootedMetricTruncation.lean} proves the actual metric axioms on an
arbitrary index type, the signatures and bounds, idempotence, full local
positive-profile invariance, marked-point offsets, coordinate differences,
and both sharpness examples. All sixteen principal axiom audits passed in
Lean~4.35.0; the file contains forty theorem declarations and uses only the
standard logical foundations. A separate exact checker examined 6,321
four-point metrics, 102,000 rooted instances, 96 large-weight adversarial
metrics, 864 full-profile comparisons, and the sharpness families for
$0\le R\le20$. The finite tests corroborate the universal proof; their
range is not the theorem's range.
